\section{The Case for the Instruments}

A cesium fountain clock will lose about one second in a hundred million years. Two laboratories on opposite sides of the world, staffed by people who share no language, politics, or religion, will read the same interferometer the same way, and when their numbers disagree past the stated uncertainty, both teams treat the disagreement as a fact to be explained rather than an opinion to be tolerated. There is nothing else in human life quite like this. Courts, parliaments, and philosophy departments should all envy it.

Set beside that record the history of the questions instruments cannot touch. Whether the world had a beginning, whether the soul survives the body, whether anything necessary stands beneath the contingent things we see---these questions are ancient, and the disputes about them are ancient too. No electrometer has ever moved when an argument for the existence of God succeeded. The debates that measurement entered, measurement largely settled; the debates it could not enter appear, from a certain distance, never to have moved at all.

From this asymmetry a conclusion has repeatedly been drawn, and it deserves to be stated at full strength, because this essay exists to argue with it. The conclusion is not that empirical method is excellent, which no one here disputes. It is that empirical method is \emph{exhaustive}: that what cannot in principle be brought before observation is not merely unsettled but unknowable, or not merely unknowable but not even a genuine question. Knowledge, on this view, has exactly one instrument panel, and whatever registers on no instrument belongs to poetry, psychology, or noise.

The pedigree of this position is distinguished. David Hume closed his \emph{Enquiry Concerning Human Understanding} with the most famous book-burning order in philosophy:

\begin{quote}
When we run over libraries, persuaded of these principles, what havoc must we make? If we take in our hand any volume; of divinity or school metaphysics, for instance; let us ask, \emph{Does it contain any abstract reasoning concerning quantity or number?} No. \emph{Does it contain any experimental reasoning concerning matter of fact and existence?} No. Commit it then to the flames: For it can contain nothing but sophistry and illusion.\endnote{David Hume, \emph{An Enquiry Concerning Human Understanding}, sect.~XII, part~III, final paragraph (E 12.34; SBN 165). Wording checked 2026-07-24 against the critical text at davidhume.org. The \emph{Enquiry} was first published in 1748; the text is in the public domain.}
\end{quote}

Note the shape of the verdict. The offending volume is not to be read and refuted; it is to be classified and destroyed. Two admissible kinds of content are named---relations of quantity, matters of experimental fact---and whatever fits neither is condemned without further hearing.

Auguste Comte gave the same verdict a philosophy of history. Each branch of knowledge, he announced, ``passes successively through three different theoretical conditions: the Theological, or fictitious; the Metaphysical, or abstract; and the Scientific, or positive.''\endnote{Auguste Comte, \emph{Cours de philosophie positive} (1830--1842), opening chapter, as condensed and translated by Harriet Martineau, \emph{The Positive Philosophy of Auguste Comte} (London: George Bell \& Sons, 1896), vol.~1, pp.~1--2. Wording checked 2026-07-24 in a scan of the 1896 Bell printing. Martineau's condensation, which Comte himself endorsed, is in the public domain.} Theology and metaphysics are not rivals to be answered but stages to be outgrown, as a man outgrows his childhood; the positive stage renounces ``the vain search after Absolute notions, the origin and destination of the universe, and the causes of phenomena'' and confines itself to laws, ``their invariable relations of succession and resemblance.'' On this account no argument against metaphysics is required, any more than an argument is required against being young.

The twentieth century sharpened the verdict into a criterion. The logical positivists of the Vienna Circle, and their most effective English advocate, A.~J. Ayer, saw that calling metaphysics \emph{false} concedes too much, since a false statement is still a statement. The stronger move was to deny that metaphysical sentences say anything at all. In \emph{Language, Truth and Logic}, published in 1936 when he was twenty-five, Ayer laid down the rule: ``The criterion which we use to test the genuineness of apparent statements of fact is the criterion of verifiability. We say that a sentence is factually significant to any given person, if, and only if, he knows how to verify the proposition which it purports to express.''\endnote{A.~J. Ayer, \emph{Language, Truth and Logic} (London: Victor Gollancz, 1936), ch.~1, ``The Elimination of Metaphysics.'' Wording checked 2026-07-24 in a scan of the Pelican reprint. Sentences with no empirical cash value may retain \emph{emotional} significance, Ayer allows, but not \emph{literal} significance. Brief quotations from this copyrighted work appear here under ordinary fair-quotation practice.} Sentences that fail the test are not mistaken; they are ``pseudo-propositions.'' And Ayer did not leave the application to the reader's imagination: ``no statement which refers to a `reality' transcending the limits of all possible sense-experience can possibly have any literal significance,'' so that ``the labours of those who have striven to describe such a reality have all been devoted to the production of nonsense.''

The consequence for religion was drawn with the same tidiness. Ayer's position was not atheism, which he considered as meaningless as theism, since the atheist denies the very sentence the theist affirms and a denial of nonsense is nonsense too. ``To say that `God exists' is to make a metaphysical utterance which cannot be either true or false.''\endnote{Ayer, \emph{Language, Truth and Logic}, ch.~6, ``Critique of Ethics and Theology.'' Ayer is explicit that his view is neither atheism nor agnosticism: the agnostic treats ``God exists'' as a genuine but undecidable proposition, whereas Ayer denies it is a proposition at all. The same chapter dissolves moral judgments into expressions of feeling---a consequence worth noticing, since it shows how much besides theology the criterion condemns.} The believer is not wrong about the world. He has failed, while appearing to speak, to say anything about the world at all.

This family of positions needs a name, and the available technical labels---positivism, verificationism, scientism---each carry historical baggage this essay does not want to adjudicate. I will call its adherents \emph{instrumentalists}: those who convert the instruments of empirical science, and the observational tests they serve, from an excellent means of knowing into the sole criterion of what can be known or what is real.\endnote{The term is coined for this essay and is defined only as stated. It is \emph{not} the doctrine called instrumentalism in the philosophy of science---the view, associated with Duhem and others, that scientific theories are calculating devices for organizing observations rather than descriptions of unobservable reality. That doctrine concerns the interpretation of science and is untouched by this essay. Nor is the target empiricism as such, the thesis that our natural knowledge originates in experience (a thesis this essay will defend in its Thomistic form), nor the methodological naturalism by which the empirical sciences rightly confine themselves to natural causes within their own work. The target is the \emph{exclusivist} conversion of empirical method into the universal test of meaning, knowledge, or reality.} The definition matters. An instrumentalist in this sense need not be a working scientist, and most working scientists are not instrumentalists in this sense; the position is a philosophical thesis about knowledge, not a laboratory practice.

One more thing should be conceded before the argument begins, because it is true and because the position's adherents have earned it. Instrumentalism is almost always adopted as a discipline of honesty. Its attraction is not contempt for meaning but hatred of wishful thinking: the suspicion that where no test is possible, assertion expands to fill the available space, and that humanity has spent millennia mistaking the echo of its own hopes for the voice of reality. That suspicion is not contemptible. It is, in fact, a moral posture---a commitment about how one \emph{ought} to form beliefs---and the irony of that fact will matter later. But a moral posture wrapped around a philosophical thesis must still submit the thesis to examination. The thesis is that only the empirically testable can be known. The examination will arrive, in due course, at an awkward question---how would one test \emph{it}?---but the question lands with its full weight only on a position seen whole. The position deserves that. It has a history, a literature, and a record of self-correction that its casual modern heirs rarely know; and the history, it turns out, is not a neutral preface to the argument but the argument's first exhibit.
